\section{Additional Experimental Details}
\label{app:recoverability_extra}

\subsection{Synthetic generation protocol}
\label{app:synthetic_protocol}
For a synthetic instance, we sample latent utilities under the hierarchical Gaussian prior
(Section~\ref{sec:model}) and construct the implied ground-truth partial orders
$\{h^{(a)}\}_{a=0}^A$.
We then generate observed traces by sampling choice sets and drawing an order from the
frontier-softmax policy with parameters $(\beta,\epsilon)$.
Unless otherwise stated, results are reported for a fixed random seed.

\paragraph{Latest run configuration.}
The synthetic results reported in the appendix are generated with $|\mathcal{A}|=8$ items and
$A=5$ assessors, each observing $N=20$ traces with choice-set sizes sampled uniformly from
$\{2,\dots,6\}$.
We fix trembling-hand noise to $\epsilon=0.05$ and sweep inverse temperature
$\beta\in\{0.8,1.5\}$ and latent poset dimension $K\in\{2,5,20\}$.
Inference uses $100{,}000$ MCMC iterations with a $50\%$ burn-in fraction, and evaluation uses
$500$ post burn-in posterior samples.

\begin{table}[t]
\centering
\caption{\textbf{Synthetic recoverability (latest run).} Cover-edge recovery metrics for assessor $a=1$
($|V|=5$). We compute F1/SHD on transitive reductions (Hasse diagrams). We report WAIC normalized per executed
step (WAIC/step) to avoid trivial scaling with dataset size. As $K$ increases, dominance becomes stricter and
the induced posets become sparser; for $K=20$, the ground-truth graph contains only a single cover edge.}
\label{tab:synthetic_latest_run}
\begin{small}
\begin{tabular}{ccccccc}
\toprule
$K$ & $\beta$ & $|E^\star|$ & F1 $\uparrow$ & SHD $\downarrow$ & ECE $\downarrow$ & WAIC/step $\downarrow$ \\
\midrule
2  & 0.8 & 4 & 0.889 & 1 & 0.090 & 0.926 \\
2  & 1.5 & 4 & 0.889 & 1 & 0.090 & 0.578 \\
5  & 0.8 & 2 & 0.667 & 1 & 0.094 & 1.099 \\
5  & 1.5 & 2 & 0.667 & 2 & 0.071 & 0.803 \\
20 & 0.8 & 1 & 1.000 & 0 & 0.018 & 1.188 \\
20 & 1.5 & 1 & 0.667 & 1 & 0.081 & 0.884 \\
\bottomrule
\end{tabular}
\end{small}
\end{table}

\subsection{WAIC normalization}
\label{app:waic_normalization}
WAIC is additive over predictive units.
If each trace $y_i$ is treated as a unit, the total WAIC scales approximately linearly with the
number of traces (and with typical trace length).
Therefore, we report WAIC normalized by the number of executed steps (WAIC/step) in the main
paper. We include unnormalized WAIC curves in the appendix for completeness.

\subsection{Identifier leakage ablation}
\label{app:id_leakage}
Workflow datasets often contain task identifiers that correlate with their position in a
workflow template.
To ensure recovery results reflect structural learning rather than ID memorization, we randomize
node IDs per workflow instance prior to inference.
We report an ablation without randomization to quantify the magnitude of this effect.
